\documentclass[11pt]{article}

\usepackage[margin=1in]{geometry}
\usepackage{amsmath,amssymb,amsthm,mathtools}
\usepackage{booktabs}
\usepackage{hyperref}
\usepackage{microtype}
\setlength{\emergencystretch}{2em}

\newtheorem{theorem}{Theorem}[section]
\newtheorem{corollary}[theorem]{Corollary}
\newtheorem{lemma}[theorem]{Lemma}
\numberwithin{equation}{section}
\theoremstyle{definition}
\newtheorem{definition}[theorem]{Definition}
\theoremstyle{remark}
\newtheorem{remark}[theorem]{Remark}

\newcommand{\C}{\mathbb C}
\newcommand{\R}{\mathbb R}
\newcommand{\K}{\mathcal K}
\newcommand{\ran}{\operatorname{ran}}
\newcommand{\kerop}{\operatorname{ker}}
\newcommand{\spec}{\operatorname{spec}}

\title{Ground-State Multiplicity Is Not a Finite Spectral Obstruction:\\
Radical Quotients for CCM Weil Matrices}
\author{}
\date{July 22, 2026}

\begin{document}
\maketitle

\begin{abstract}
The finite spectral construction of Connes--Consani--\allowbreak Moscovici starts from a
positive-semidefinite
shift of a truncated Weil matrix and assumes that its radical is generated by one even vector.
We prove an arbitrary-radical replacement.  Every reflection-graded divided-difference matrix in
their class admits a rank-one boundary completion which is self-adjoint for the semidefinite Weil
metric and preserves its entire radical.  After quotienting by that radical, the completed clock is
an ordinary self-adjoint operator.  In Gram coordinates the quotient is the explicit Hermitian
matrix
\[
                         C=BAR,
\]
where $T=B^*B$ is any full-row-rank Gram factor and $BR=I$.  Consequently its characteristic
polynomial has only real zeros, without simplicity or parity assumptions on the ground space.

This does not prove that the original scalar ground-vector Fourier transform used by
Connes--Consani--Moscovici survives a multiple ground state.  It produces a different, canonical
object: the full-radical quotient determinant.  We show that locally uniform convergence of
normalized quotient determinants to Riemann's $\Xi$-function implies the Riemann hypothesis.
Thus the finite simple-even obstruction is replaced by one sharply stated analytic problem:
divisor-preserving convergence of radical-quotient determinants.  The general theory of operators
on semidefinite inner-product spaces is classical; the contribution is the boundary completion and
its application to the CCM divided-difference system.
\end{abstract}

\section{The obstruction in the CCM construction}

Let $Q=Q_{\lambda,N}$ be a finite truncated Weil matrix in the Fourier basis indexed by
$-N,\ldots,N$.  Connes--Consani--Moscovici show that it has the form
\begin{equation}
 Q_{ii}=a_i,\qquad
 Q_{ij}=\frac{b_i-b_j}{i-j}\quad(i\ne j),
\end{equation}
with $a_{-i}=a_i$ and $b_{-i}=-b_i$.  If
\[
 D e_i=i e_i,\qquad \Gamma e_i=e_{-i},\qquad
 \eta=\sum_i e_i,\qquad \beta=\sum_i b_i e_i,
\]
then
\begin{equation}
 DQ-QD=|\beta\rangle\langle\eta|-|\eta\rangle\langle\beta|,
 \qquad Q\Gamma=\Gamma Q.
\end{equation}
The vector $\eta$ is even and $\beta$ is odd.

Put $\mu=\min\spec Q$ and $T=Q-\mu I$.  Then $T\ge0$, but its radical
$\K=\kerop T$ need not be one-dimensional.  The published construction assumes
\begin{equation}
 \K=\C\xi,\qquad \Gamma\xi=\xi,\qquad
 \langle\eta,\xi\rangle\ne0.
\end{equation}
That assumption selects the correction
\[
 D'=D-|D\xi\rangle\langle\eta|.
\]
The central question here is different: is condition (1.3) needed to obtain a finite
self-adjoint spectral object, or only to identify that object with one scalar Fourier transform?

\section{Boundary-channel completion of the full radical}

We first isolate the finite algebra behind (1.2).  All spaces in this section are finite
dimensional real or complex Hilbert spaces.

\begin{theorem}[Boundary-channel trichotomy]
\label{thm:completion}
Let $T,D,\Gamma$ be self-adjoint, let $T\Gamma=\Gamma T$, and let
\[
 \Gamma\eta=\eta,\qquad \Gamma\beta=-\beta,
 \qquad DT-TD=|\beta\rangle\langle\eta|-|\eta\rangle\langle\beta|.
\]
Writing $\K=\kerop T$, one of the following three constructions always applies.
\begin{enumerate}
\item If $\eta|_{\K}\ne0$, there is $e\in\K$ with
      $\langle\eta,e\rangle=1$ and $\langle\beta,e\rangle=0$, and
      \[
      A_E=D-|De\rangle\langle\eta|.
      \]
\item If $\eta|_{\K}=0$ and $\beta|_{\K}\ne0$, there is $o\in\K$ with
      $\langle\beta,o\rangle=1$ and $\langle\eta,o\rangle=0$, and
      \[
      A_O=D-|Do\rangle\langle\beta|.
      \]
\item If $\eta|_{\K}=\beta|_{\K}=0$, there is $y$ with $Ty=\beta$, and
      \[
      A_N=D+|y\rangle\langle\eta|.
      \]
\end{enumerate}
In every case the resulting operator $A$ satisfies
\begin{equation}
                         TA=A^*T,
              \qquad A\K\subseteq\K.
\end{equation}
The first two active boundary channels are mutually exclusive.
\end{theorem}

\begin{proof}
For $k\in\K$, (2.1) gives the radical-transport identity
\begin{equation}
 TDk=-\beta\langle\eta,k\rangle
       +\eta\langle\beta,k\rangle.
\end{equation}
Suppose first that $\eta|_{\K}\ne0$.  Reflection averaging selects an even $e\in\K$ with
$\langle\eta,e\rangle=1$; parity gives $\langle\beta,e\rangle=0$.  Equation (2.6) yields
$TDe=-\beta$.  Symmetry then forces $\beta\perp\K$.  Substitution in (2.2) gives
\[
 TA_E=TD+|\beta\rangle\langle\eta|
     =DT+|\eta\rangle\langle\beta|=A_E^*T,
\]
and (2.6) shows that $A_E\K\subseteq\K$.  The odd case is the same argument with
$\eta$ and $\beta$ interchanged.

In the neutral case both boundary functionals annihilate $\K$.  Since $T=T^*$,
$(\kerop T)^\perp=\ran T$, so $\beta=Ty$ for some $y$.  Direct substitution of (2.4) into
(2.1) proves (2.5).  Finally, an active even vector makes $\beta$ vanish on all of $\K$, while an
active odd vector has nonzero $\beta$-value, proving mutual exclusion.
\end{proof}

Theorem~\ref{thm:completion} uses neither positivity nor a preferred ground vector.  Positivity
enters only when the radical is converted into an honest Hilbert quotient.

\section{The radical-quotient Hermitian law}

The general fact that a bounded self-adjoint operator on a semidefinite inner-product space
descends to a self-adjoint operator on the quotient by the isotropic part is classical; see
\cite{Bovdi2020}.  The next formula is the coordinate form needed for the CCM system.

\begin{theorem}[Gram-coordinate radical quotient]
\label{thm:gram}
Let $B:\C^d\to\C^r$ be surjective, let $R:\C^r\to\C^d$ satisfy $BR=I_r$, and set
$T=B^*B$.  If $A:\C^d\to\C^d$ satisfies
\begin{equation}
                         TA=A^*T,
\end{equation}
then
\begin{equation}
                         C=BAR
\end{equation}
is Hermitian.  If in addition $A\kerop B\subseteq\kerop B$, then
\begin{equation}
                         CB=BA.
\end{equation}
Thus $C$ represents the operator induced by $A$ on $\C^d/\kerop T$.
\end{theorem}

\begin{proof}
Multiplying (3.1) on the left by $R^*$ and on the right by $R$ gives
\[
 BAR=R^*B^*BAR=R^*A^*B^*BR=R^*A^*B^*=C^*.
\]
For (3.3), the operator $I-RB$ has image in $\kerop B$.  Radical preservation therefore gives
$BA(I-RB)=0$, or $BA=BARB=CB$.  Finally,
$\kerop T=\kerop(B^*B)=\kerop B$, so (3.3) is precisely the quotient intertwining law.
\end{proof}

\begin{corollary}[Arbitrary-radical CCM real spectrum]
\label{cor:ccm}
Let $T$ be any positive-semidefinite scalar shift of a finite reflection-graded
divided-difference matrix (1.1).  Choose a full-row-rank Gram factor $T=B^*B$ and a right inverse
$R$ of $B$.  Then one of the boundary completions in Theorem~\ref{thm:completion} produces an
operator $A$ for which $C=BAR$ is Hermitian and $CB=BA$.  In particular,
\begin{equation}
                         \det(C-zI)
\end{equation}
has only real zeros.  No dimension or parity condition on $\kerop T$ is required.
\end{corollary}

\begin{proof}
Every finite positive-semidefinite matrix has a full-row-rank Gram factor whose row number is its
rank.  Apply Theorem~\ref{thm:completion}, followed by Theorem~\ref{thm:gram}.  Hermitian matrices
have real spectrum, which proves the last assertion.
\end{proof}

Corollary~\ref{cor:ccm} is not the assertion that every ground vector has a real-rooted Fourier
transform.  When the radical has dimension greater than one, the correct object produced by the
argument is the determinant of the complete radical quotient.  This distinction is essential:
the metric contains no spectral information on the residual radical itself.

\section{The quotient-determinant route to RH}

Let $C_j$ be radical-quotient matrices obtained from a cofinal family of CCM truncations.  Let
$u_j$ be entire and nowhere zero, and define
\begin{equation}
                       F_j(z)=u_j(z)\det(C_j-zI).
\end{equation}

\begin{theorem}[Radical-quotient determinant criterion]
\label{thm:rh}
If $F_j\to\Xi$ locally uniformly on $\C$, then the Riemann hypothesis holds.
\end{theorem}

\begin{proof}
Every zero of $F_j$ is real by Corollary~\ref{cor:ccm}, because $u_j$ has no zeros.  Suppose
$\Xi(z_0)=0$ with $\operatorname{Im}z_0\ne0$.  Choose a closed disk about $z_0$ disjoint from the
real axis and containing no zero of $\Xi$ on its boundary.  Hurwitz's theorem, equivalently the
argument principle under uniform boundary convergence, forces $F_j$ to have a zero in that disk
for all sufficiently large $j$, a contradiction.  Hence every zero of $\Xi$ is real, which is RH
in the centered coordinate.
\end{proof}

Theorem~\ref{thm:rh} is an endpoint, not a disguised proof of convergence.  It changes the CCM
research problem in a concrete way:
\[
\begin{array}{c}
\text{simple even ground vector}\quad+\quad
\text{scalar Fourier-transform convergence}
\\[2mm]
\Downarrow\quad\text{is replaced by}\quad\Downarrow
\\[2mm]
\text{full-radical boundary completion}\quad+\quad
\text{quotient-determinant convergence}.
\end{array}
\]
The first item on the second line is unconditional at every finite truncation.  The second is the
open analytic core.

\section{What has and has not been removed}

The result removes ground-state simplicity and parity from the finite real-spectrum mechanism.
It does not prove the two assertions below.

\begin{enumerate}
\item It does not identify (3.4), in a multiple-radical regime, with the Fourier transform of one
      chosen ground vector.  That identification is special to a one-dimensional active channel.
\item It does not prove that normalized quotient determinants converge to $\Xi$.  A successful
      proof must control the divisor after the entire residual radical has been cancelled, not
      merely show convergence of matrix entries or of a selected eigenvector.
\end{enumerate}

This gives two immediate falsification tests for the route.  First, along any proposed cofinal
regime the quotient determinant must be stable under the radical cancellation and admit a
nowhere-zero normalization.  Second, its zero counting measures must be tight on every compact
set and converge to the divisor of $\Xi$.  Failure of either property kills the route without any
appeal to RH itself.

The universal quotient step is not claimed as new: Lemma~3 of \cite{Bovdi2020} gives the
semiunitary-space principle in a substantially broader canonical-form setting.  The candidate new
content is the exhaustive rank-one boundary completion forced by the CCM displacement law and the
resulting arbitrary-radical determinant route.  The papers \cite{CCM2025,CvS2025} use a simple even
ground state and do not state this replacement.

\begin{thebibliography}{9}

\bibitem{CCM2025}
A.~Connes, C.~Consani, and H.~Moscovici,
\emph{Zeta Spectral Triples}, arXiv:2511.22755 (2025),
\url{https://arxiv.org/abs/2511.22755}.

\bibitem{CvS2025}
A.~Connes and W.~D. van Suijlekom,
\emph{Quadratic Forms, Real Zeros and Echoes of the Spectral Action},
arXiv:2511.23257 (2025),
\url{https://arxiv.org/abs/2511.23257}.

\bibitem{Bovdi2020}
V.~A. Bovdi, T.~Klymchuk, T.~Rybalkina, M.~A. Salim, and V.~V. Sergeichuk,
\emph{Operators on positive semidefinite inner product spaces},
Linear Algebra Appl. \textbf{596} (2020), 82--105,
\url{https://arxiv.org/abs/2003.06919}.

\end{thebibliography}

\end{document}
